Species in a plasma exchange their mass, momentum, and energy primarily through local collisions. The collisions described in this section are assumed to be entirely binary; i.e.\ involving two particles at a
time. There are two interactions that can follow a local collision of unlike particles, namely reaction and relaxation. Reactions convert a neutral particles into ions or vice versa, and relaxations drive the
system towards equilibrium in absence of external sources.

\subsection{Relaxations}
\label{sect:relaxations}
A collision between a charged particle $j$ and a neutral particle results in a force roughly equal to
\begin{equation}
    \mathbf{f}_{jn} = \frac{m_j n_j (\mathbf{u}_j - \mathbf{u})}{\tau_{jn}}.
    \label{eqn:momentum_relaxation_jn}
\end{equation}
Here the double subscript $jn$ refers to a transfer from a particle $j$ to a neutral particle. By Newton's third law of motion, $\mathbf{f}_{jn} = -\mathbf{f}_{nj}$. The mean time between collisions $\tau_{jn}$
is approximately constant and is the reciprocal of the momentum-transfer collision frequency
\begin{equation}
    \nu_{jn}^M \equiv \frac{1}{\tau_{jn}} = n \sigma^{M}_{jn} \lVert \mathbf{u}_j - \mathbf{u} \rVert,
    \label{eqn:nu_jn_momentum}
\end{equation}
with $\sigma^{M}_{jn}$ denotes the equivalent cross section of momentum transfer between $j$ and $n$. It is typically a function of charged particle speed; the average value over all velocities in a Maxwellian
distribution is typically used in Equation \ref{eqn:nu_jn_momentum}.

Collisions between two charged particles are Coulomb collisions mediated by the electric field. Like-particle (electron-electron or ion-ion) collisions are often neglected because they preserve the momentum and
energy between themselves, and they contribute very little to diffusion. This is in contrast with electron-ion collisions, where there exists a net transfer of momentum (and energy) between the two particles. In a
similar formulation as Equation \ref{eqn:momentum_relaxation_jn}, the force exerted between an ion and an electron during their collision is
\begin{equation}
    \mathbf{f}_{ei} = m_e n_e (\mathbf{u}_i - \mathbf{u}_e) \nu^M_{ei},
    \label{eqn:momentum_relaxation_ei}
\end{equation}
and the electron-ion momentum-transfer collision frequency $\nu^M_{ei}$ is
\begin{equation}
    \nu^M_{ei} = C \frac{n_i Z^2 e^4}{(4\pi\varepsilon_0)^2 \sqrt{m_e (k_B T_e)^3}} \ln\Lambda_e,
    \label{eqn:nu_ei_momentum}
\end{equation}
where $C$ is a prefactor on the order of $\mathcal{O}(1)$ that differs depending on the text. Chen gives $C = \pi$, while the FLASH code uses $C = \frac{4}{3} \sqrt{2\pi}$ \cite{Chen}\cite{FLASH}. The latter will be
used in this model. The Coulomb logarithm $\ln\Lambda$ is a correction factor that accounts for the cumulative effect of many small-angle deflections in a collision and is given by
\begin{equation}
    \ln\Lambda_e \equiv \ln \left[ \frac{4\pi (\varepsilon_0 k_B T_e)^{3/2}}{Ze^3 \sqrt{n_e}} \right].
    \label{eqn:ln_lambda}
\end{equation}
In both Equations \ref{eqn:nu_ei_momentum} and \ref{eqn:ln_lambda}, the average ionization number $Z$ can be assumed to be 1.

For all types of collisions, the energy loss due to relaxation term can be expressed as a heat source that flows from particle $j$ to $j'$
\begin{equation}
    q'''_{jj'} = -q'''_{j'j} = \frac{3}{2} n_j \nu_{jj'}^K k_B (T_j - T_{j'}).
    \label{eqn:energy_relaxation}
\end{equation}

The energy-transfer collision frequency $\nu^K$ is related to its momentum counterpart $\nu^M$ by
\begin{equation}
    \frac{\nu_{jj'}^K}{\nu_{jj'}^M} = \frac{2m_j}{m_j + m_{j'}}.
    \label{eqn:collision_frequency_ratio}
\end{equation}

The index $j$ in Equation \ref{eqn:collision_frequency_ratio} refers to the lighter particle out of the pair, so it is always the electron if it is involved in a collision, and the ion in a ion-neutral collision.
In the former case, $\nu_{ej'}^K \ll \nu_{ej'}^M$, which means energy is relaxed at a much slower rate than momentum is. Physically, upon colliding with a heavy particle, an electron is frequently scattered before
a significant portion of its kinetic energy is lost. In an ion-neutral collision, due to the near-identical mass of the two particles, energy and momentum are relaxed on the same timescale -- i.e., $\nu_{in}^K = \nu_{in}^M$.